\section{Scope, evidence levels, and adjacent literature}

\subsection{Epistemic firewall}

We distinguish four levels of statement throughout:

\begin{enumerate}
  \item \textbf{Operational:} specified coherent evolutions yield an interference signal and controls do not.
  \item \textbf{Model-relative historical:} components of a declared decomposition carry different internal records at intermediate times.
  \item \textbf{Interpretive:} one may describe those components using a preferred interpretation of quantum mechanics.
  \item \textbf{Unsupported:} the experiment accesses a pre-existing macroscopic alternative world or establishes a unique ontology.
\end{enumerate}

Only the first level is directly measured. The second follows from a validated model and declared tensor decomposition. The third is underdetermined by shared unitary predictions. The fourth is outside the protocol.

\subsection{Known ingredients}

Relative-state language and decoherence motivate discussions of correlated components and records \cite{everett1957relative,zurek2003decoherence}. Consistent-histories methods describe sequences and their decoherence functionals \cite{griffiths1984consistent,paz1993environment}, although alternatives that are deliberately recombined are not simultaneously a classical consistent family. Reversible computation supplies compute--copy-or-phase--uncompute constructions \cite{bennett1973logical}. Wave--particle duality and quantum erasure quantify the loss and conditional recovery of interference when path information exists \cite{englert1996fringe,kim2000delayed}. Measurement reversal and error-correction demonstrations establish that selected measurement interactions can be undone under controlled conditions \cite{katz2008reversal,schindler2013undoing}.

For multitime structure, quantum combs and process tensors treat an experiment as a process with intervention slots rather than as a single endpoint channel \cite{chiribella2009networks,pollock2018processes,hoffreumon2021multiround}. Quantum causal discovery makes interventions central to the distinction between correlation and causation \cite{giarmatzi2018causal}. Temporal studies can diagnose instrument-specific memory effects, bound effective environment dimension, and, under stronger assumptions, certify a quantum memory without trusting its implementation \cite{guo2021memory,vieira2024dimension,sekatski2023memory}. A January 2026 preprint reports a trapped-ion proof of principle for device-independent two-time memory certification \cite{santos2026deviceindependent}; its peer-review status at our cutoff is provisional.

Wigner-friend-inspired protocols show how a small physical record can play the functional role of an internal observer, including explicit analyses of memory erasure \cite{frauchiger2018quantum,proietti2019observer,bong2020strong,elouard2021erasing}. These works do not place a conscious human observer in coherent superposition. The memory-erasure result is a particularly direct antecedent and prevents a novelty claim based merely on a friend acquiring and then losing a record.

Quantum zero-knowledge proof systems establish rigorous ways for a verifier to learn validity without learning a witness, including QMA, classical-verifier protocols for quantum computations, non-destructive state proofs, and superposition-secure transcript models \cite{broadbent2016qma,vidick2020classical,colisson2025nondestructive,coladangelo2026mpc}. Their security definitions are computational or information-theoretic statements about simulators and adversaries. Our zero-record criterion instead asks whether physical residual systems remain correlated with distinctions inside a predicate class. The analogy is generative but not an equivalence.

Quantum algorithmic measurements (QUALMs) provide a complexity-theoretic language for experimental access and exhibit tasks with exponential separations between coherent and incoherent access \cite{aharonov2022qualm}. Computational mechanics supplies irreducible predictive-memory measures, and quantum models or adaptive agents can reduce memory costs for selected stochastic processes \cite{gu2012complexity,elliott2022agents}. These results suggest the right form of a future originality claim: not that a circuit used superposition, but that a precisely defined historical-certification task has a provable resource separation.

\subsection{Evidence and novelty discipline}

We classify claims as known results, derivations, numerical observations, conjectures, or open objectives. The accompanying literature review is a scoping review rather than a PRISMA systematic review. Several directly adjacent manuscripts appeared as 2026 preprints and remain provisional at the cutoff date. Consequently, the phrase ``not found in this search'' is never promoted to ``never proposed.'' A priority claim requires database-complete searches, forward and backward citation tracing, and specialist review.

The elementary history--predicate--uncompute circuit is not new. Generic memory--visibility inequalities are also heavily constrained by existing complementarity and channel-recovery results \cite{englert1996fringe,kretschmann2008information}. The candidate open intersection is the joint task of causal-memory certification, within-predicate zero-record leakage, and subsequent recoherence under an explicit access model.

\subsection{Adversarial novelty map}

Table~\ref{tab:novelty} records which intuitive headlines are already occupied. It is deliberately stricter than a conventional motivation section.

\begin{table}[ht]
\centering
\caption{Adjacency audit at the 12 August 2026 cutoff. The final row is a candidate research task, not an established priority claim.}
\label{tab:novelty}
\small
\begin{tabularx}{\textwidth}{@{}>{\raggedright\arraybackslash}p{0.23\textwidth}Y Y@{}}
\toprule
Candidate headline & Existing territory & Remaining requirement \\
\midrule
Compute a predicate and erase work & Reversible computation and phase kickback & No novelty without a new access or resource theorem \\
Erase an internal observer's memory & Wigner-friend memory erasure is explicit prior work & Certify causal memory and nonvacuous final privacy \\
Bound memory by temporal data & Process-tensor and memory-dimension witnesses exist & Combine the witness with a later coherent inverse \\
Hide a quantum witness or transcript & Quantum zero-knowledge protocols cover strong adversarial models & Relate simulator security to physical final-record decoupling \\
Show coherent experimental advantage & QUALMs already give exponential separations for other tasks & Prove a separation for the zero-record task itself \\
Trade objectivity against recovery & Darwinism, encoding transitions, and QEC tradeoffs are active & Only a narrower operational theorem may remain \\
Causal memory + predicate-only disclosure + final decoupling + recoherence & No single primary source located by the bounded search & Formal anti-shortcut model and nontrivial theorem still required \\
\bottomrule
\end{tabularx}
\end{table}

\subsection{Direct 2026 priority risks}

The current literature makes interpretation-led framing especially risky. Three January 2026 preprints use inter-branch communication language, compile Wigner-friend-style message circuits, or benchmark them on superconducting processors \cite{violaris2026branches,altman2026wigner,cogburn2026interbranch}. Another January preprint studies approximate records, decoherence, recoherence, and recovery in long isolated histories \cite{strasberg2026records}. These manuscripts are not treated as settled peer-reviewed results, but they are direct priority evidence against presenting branch communication, memory erasure, or generic recoherence as new.

Quantum Darwinism formalises redundant proliferation of classical records, and models already connect redundant objective information with consistent histories and quantum encoding \cite{blumekohout2006darwinism,riedel2016objective,ferte2024transition}. A March 2026 peer-reviewed model further distinguishes approximate decoherent histories from the redundant records associated with objectivity \cite{ferte2026histories}. A preprint posted eight days before our cutoff claims an exact QEC--Darwinism tradeoff and a model-independent no-go statement \cite{maity2026tradeoff}. Its generality requires independent scrutiny, but its existence makes objectivity versus recoverability a poor novelty headline.

Thermodynamic erasure is also not an empty intersection. Complexity-constrained quantum thermodynamics derives work--complexity tradeoffs for reset \cite{munson2025complexity}, while microscopic measurement models study records, reset, finite-time control, and work \cite{latune2025measurement}. Our inverse is logically reversible; assigning \(k_BT\ln2\) to every uncomputed bit would conflate uncomputation with blind erasure.

\subsection{Macroscopic scale and nonunitarity}

Matter-wave experiments have extended coherent interference to increasingly massive molecules and nanoparticles, including a 2026 sodium-nanoparticle result \cite{pedalino2026nanoparticle}. This calibrates centre-of-mass coherence, not reversible control of a macroscopic internal observer and its environment. Objective-collapse models make parameterised predictions and are experimentally constrained, including by a March 2026 XENONnT spontaneous-radiation search \cite{bassi2013collapse,nimmrichter2014optomechanical,arnquist2022majorana,aprile2026collapse}. An unexplained visibility loss in the present protocol would first indicate an incomplete noise or control model; it would not by itself identify fundamental nonunitarity.

\subsection{Conservative novelty statement}

The dated search found primary literature for every clause separately but did not locate a protocol that jointly imposes: (i) an interventional lower bound on causal memory, (ii) disclosure limited to a non-injective predicate, (iii) final decoupling of every other accessible history record, (iv) recovered interference, and (v) an explicit shortcut class. This is a bounded negative search result, not proof of absence. A first submission must add citation-network traversal, subscription-database searches, patent and thesis checks, and independent specialist review before using a phrase such as \emph{to our knowledge}.
